\Askhsh{15012}{2}
\begin{ylh}{1}
\item 4.1 Πολυώνυμα
\item 4.2 Διαίρεση πολυωνύμων
\end{ylh}
Η διαίρεση ενός πολυωνύμου $P(x)$ με το $x - 3$ έχει πηλίκο $x^2 + 2$ και υπόλοιπο 4.
\begin{alist}
\item Να γράψετε την ταυτότητα της παραπάνω διαίρεσης.
\monades{8}
\item Να δείξετε ότι $P(x) = x^3 - 3x^2 + 2x - 2$.
\monades{8}
\item Είναι το $x = 3$ ρίζα του πολυωνύμου $P(x)$; Να αιτιολογήσετε την απάντησή σας.
\monades{9}
\end{alist}

